\documentclass[11pt]{article}
\usepackage{amsmath,amssymb,amsthm}
\usepackage[margin=1in]{geometry}
\usepackage{hyperref}

\newtheorem{lemma}{Lemma}
\newtheorem{theorem}{Theorem}
\newtheorem{corollary}{Corollary}
\newtheorem{remark}{Remark}

\title{Doc 07 (Referee-proof): A Uniform Wilson Hessian Bound for SU(3) in Hilbert--Schmidt Norm}
\author{Project extraction / consolidation}
\date{\today}

\begin{document}
\maketitle

\section{Purpose}

This note isolates a single analytic statement that is both:
(i) genuinely \emph{global} (independent of the field amplitude $A$ and the lattice volume), and
(ii) directly useful for turning numerical convexity scans into \emph{validated} convexity certificates.

The key technical point is to work in exponential coordinates $U_\ell = \exp(A_\ell)$ with $A_\ell\in\mathfrak{su}(3)$ and to use the Hilbert--Schmidt (HS) norm, which is invariant under left/right multiplication by unitary matrices. This removes spurious factors like $e^{c\|A\|}$ that appear if one treats $\exp(A)$ as an arbitrary matrix.

\section{Setup}

Let $G=\mathrm{SU}(3)$ and $\mathfrak{g}=\mathfrak{su}(3)$ (skew-Hermitian, traceless $3\times 3$ matrices).
For $X\in\mathbb{C}^{3\times 3}$ define the HS norm
\[
\|X\|_{HS} := \sqrt{\mathrm{Tr}(X^\dagger X)}.
\]
For a finite lattice with oriented link set $\mathcal{E}$, define the configuration space
\[
\mathcal{X} := \mathfrak{g}^{\mathcal{E}}
\]
with inner product and norm
\[
\langle A,B\rangle_{\mathcal{X}} := \sum_{\ell\in\mathcal{E}} \mathrm{Tr}(A_\ell^\dagger B_\ell),
\qquad
\|A\|_{\mathcal{X}}^2 := \sum_{\ell\in\mathcal{E}} \|A_\ell\|_{HS}^2.
\]

For each oriented link $\ell\in\mathcal{E}$ write
\[
U_\ell(A) := e^{A_\ell}\in SU(3).
\]
If $\bar\ell$ denotes the reversed orientation, set $U_{\bar\ell}(A) := U_\ell(A)^{-1} = U_\ell(A)^\dagger = e^{-A_\ell}$.

For an oriented plaquette $p=(\ell_1,\ell_2,\ell_3,\ell_4)$ define the plaquette matrix
\[
U_p(A) := U_{\ell_1}(A)U_{\ell_2}(A)U_{\ell_3}(A)U_{\ell_4}(A)\in SU(3),
\]
and the Wilson plaquette action
\[
S_W(A) := \beta \sum_{p\in\mathcal{P}} \left(1-\frac{1}{3}\mathrm{Re\,Tr}\,U_p(A)\right).
\]
The constant term $\beta\sum_p 1$ plays no role for derivatives and will be ignored when convenient.

We view $\nabla^2 S_W(A)$ as a self-adjoint operator on $\mathcal{X}$, and denote its operator norm by
\[
\|\nabla^2 S_W(A)\|_{\mathrm{op}}
:=
\sup_{\substack{H\neq 0\\K\neq 0}}
\frac{|\langle H,\nabla^2 S_W(A)K\rangle_{\mathcal{X}}|}{\|H\|_{\mathcal{X}}\|K\|_{\mathcal{X}}}.
\]

\section{Uniform bounds for derivatives of the matrix exponential}

\begin{lemma}[First and second Fr\'echet derivatives of $\exp$ in HS norm]
\label{lem:exp-derivatives}
Let $A\in\mathfrak{su}(3)$ and $H,K\in\mathfrak{su}(3)$. Then the Fr\'echet derivatives satisfy
\[
De^{A}[H]=\int_0^1 e^{(1-s)A}\,H\,e^{sA}\,ds,
\]
\[
D^2e^{A}[H,K]=\int_0^1\!\!\int_0^s e^{(1-s)A}\,H\,e^{(s-u)A}\,K\,e^{uA}\,du\,ds
\;+\;
(H\leftrightarrow K),
\]
and, crucially, the HS bounds
\[
\boxed{\ \|De^A[H]\|_{HS}\le \|H\|_{HS}\ },
\qquad
\boxed{\ \|D^2e^A[H,K]\|_{HS}\le \|H\|_{HS}\|K\|_{HS}\ }.
\]
The same bounds hold for $e^{-A}$.
\end{lemma}

\begin{proof}
Since $A\in\mathfrak{su}(3)$, $e^{tA}$ is unitary for all $t\in\mathbb{R}$. The HS norm is unitarily invariant:
$\|UXV\|_{HS}=\|X\|_{HS}$ for all unitary $U,V$.
Therefore for each $s$,
$\|e^{(1-s)A}He^{sA}\|_{HS}=\|H\|_{HS}$, and integrating gives $\|De^A[H]\|_{HS}\le\|H\|_{HS}$.

For the second derivative, by submultiplicativity in operator$\times$HS norms and $\|M\|_{\mathrm{op}}\le\|M\|_{HS}$,
\[
\|e^{(1-s)A}H e^{(s-u)A}K e^{uA}\|_{HS}
=\|H e^{(s-u)A}K\|_{HS}
\le \|H\|_{\mathrm{op}}\|K\|_{HS}
\le \|H\|_{HS}\|K\|_{HS}.
\]
The triangle-domain integral has area $1/2$, and there are two symmetric terms, so the stated bound with constant $1$ follows.
The same reasoning applies to $e^{-A}$ since $e^{-tA}$ is also unitary.
\end{proof}

\section{Derivative bounds for the plaquette trace}

Define the basic plaquette trace functional on $(SU(3))^4$:
\[
\phi(X_1,X_2,X_3,X_4)
:=
-\frac{1}{3}\mathrm{Re\,Tr}(X_1X_2X_3X_4).
\]
For each $i$ denote by $D_i\phi$ and $D^2_{ij}\phi$ the first and second Fr\'echet derivatives in the $X_i$ variables.

\begin{lemma}[Uniform first and mixed second derivative bounds for $\phi$]
\label{lem:trace-derivatives}
For all $X_1,\dots,X_4\in SU(3)$ and all variations $\Delta_i,\Delta_j\in\mathbb{C}^{3\times 3}$,
\[
|D_i\phi[\Delta_i]|\le \frac{1}{\sqrt{3}}\|\Delta_i\|_{HS},
\]
and for $i\neq j$,
\[
|D^2_{ij}\phi[\Delta_i,\Delta_j]|\le \frac{1}{3}\|\Delta_i\|_{HS}\,\|\Delta_j\|_{HS}.
\]
Moreover $D^2_{ii}\phi\equiv 0$ (linearity in each $X_i$ separately).
\end{lemma}

\begin{proof}
Fix $i$. By product rule,
$D_i\phi[\Delta_i]=-(1/3)\mathrm{Re\,Tr}(U\Delta_i V)$ for some unitary $U,V$ (products of the other $X_k$). Then
\[
|D_i\phi[\Delta_i]|
\le \frac{1}{3}|\mathrm{Tr}(U\Delta_i V)|
= \frac{1}{3}|\mathrm{Tr}(\Delta_i W)|,\quad W:=VU\in SU(3).
\]
Using HS Cauchy--Schwarz, $|\mathrm{Tr}(\Delta_i W)|\le \|\Delta_i\|_{HS}\|W\|_{HS}$, and $\|W\|_{HS}=\sqrt{\mathrm{Tr}(I)}=\sqrt{3}$, giving the first bound.

For $i\neq j$, $D^2_{ij}\phi[\Delta_i,\Delta_j]=-(1/3)\mathrm{Re\,Tr}(U\Delta_i V\Delta_j W)$ with $U,V,W$ unitary. Using $\mathrm{Re}(\cdot)\le|\cdot|$, cyclicity, and the Frobenius inequality $|\mathrm{Tr}(AB)|\le \|A\|_{HS}\|B\|_{HS}$,
\[
|\mathrm{Tr}(U\Delta_i V\Delta_j W)|
=|\mathrm{Tr}((U\Delta_i V)(\Delta_j W))|
\le \|U\Delta_i V\|_{HS}\,\|\Delta_j W\|_{HS}
=\|\Delta_i\|_{HS}\|\Delta_j\|_{HS}.
\]
The $D^2_{ii}\phi=0$ statement follows because $\phi$ is linear in each argument separately.
\end{proof}

\section{Uniform Hessian bound for the Wilson action}

Let $p=(\ell_1,\ell_2,\ell_3,\ell_4)$ be a plaquette and define $S_p(A):=\beta\,\phi(U_{\ell_1}(A),\dots,U_{\ell_4}(A))$.
The total Wilson action is $S_W(A)=\sum_{p\in\mathcal{P}}S_p(A)$.

\begin{lemma}[Per-plaquette block bounds]
\label{lem:plaquette-blocks}
Fix a plaquette $p$ with links $\ell_1,\dots,\ell_4$.
Then the Hessian of $S_p$ as an operator on $\mathfrak{su}(3)^4$ has block operator norms bounded by:
\[
\|(\nabla^2 S_p(A))_{ij}\|_{\mathrm{op}}
\le
\begin{cases}
\beta/\sqrt{3}, & i=j,\\[3pt]
\beta/3, & i\neq j,
\end{cases}
\qquad\text{for all }A.
\]
\end{lemma}

\begin{proof}
Write $X_i(A)=U_{\ell_i}(A)=e^{\pm A_{\ell_i}}$ depending on orientation; by Lemma~\ref{lem:exp-derivatives} the sign is irrelevant for the norms.

For $i\neq j$, $D^2_{A_i,A_j}S_p$ involves only the mixed derivative of $\phi$ composed with first derivatives of $X_i$ and $X_j$:
\[
D^2_{A_i,A_j}S_p[H_i,K_j]
=
\beta\,D^2_{ij}\phi\big[DX_i[H_i],DX_j[K_j]\big].
\]
By Lemmas~\ref{lem:trace-derivatives} and \ref{lem:exp-derivatives},
\[
|D^2_{A_i,A_j}S_p[H_i,K_j]|
\le \beta\cdot \frac{1}{3}\|DX_i[H_i]\|_{HS}\|DX_j[K_j]\|_{HS}
\le \beta\cdot \frac{1}{3}\|H_i\|_{HS}\|K_j\|_{HS}.
\]
By duality this implies the block operator norm bound $\|(\nabla^2 S_p)_{ij}\|_{\mathrm{op}}\le \beta/3$.

For $i=j$, $D^2_{ii}\phi\equiv 0$, so the only second derivative contribution comes from the second derivative of the exponential:
\[
D^2_{A_i,A_i}S_p[H_i,K_i]
=
\beta\,D_i\phi\big[D^2X_i[H_i,K_i]\big].
\]
Using Lemmas~\ref{lem:trace-derivatives} and \ref{lem:exp-derivatives},
\[
|D^2_{A_i,A_i}S_p[H_i,K_i]|
\le \beta\cdot\frac{1}{\sqrt{3}}\|D^2X_i[H_i,K_i]\|_{HS}
\le \beta\cdot\frac{1}{\sqrt{3}}\|H_i\|_{HS}\|K_i\|_{HS},
\]
so $\|(\nabla^2 S_p)_{ii}\|_{\mathrm{op}}\le \beta/\sqrt{3}$.
\end{proof}

\begin{theorem}[Uniform (field-independent) Wilson Hessian bound in 4D]
\label{thm:uniform-wilson-hessian}
On a 4D hypercubic lattice, each oriented link belongs to at most $n_p=6$ plaquettes.
For all $A\in\mathcal{X}$,
\[
\boxed{
\|\nabla^2 S_W(A)\|_{\mathrm{op}}
\le
C_W\,\beta,
\qquad
C_W := n_p\left(1+\frac{1}{\sqrt{3}}\right)=6\left(1+\frac{1}{\sqrt{3}}\right).
}
\]
In particular, a simpler but looser bound is $\|\nabla^2 S_W(A)\|_{\mathrm{op}}\le 8\sqrt{3}\,\beta$.
\end{theorem}

\begin{proof}
View $\nabla^2 S_W(A)$ as a block matrix indexed by links:
\[
\nabla^2 S_W(A) = \big(H_{\ell\ell'}(A)\big)_{\ell,\ell'\in\mathcal{E}},
\quad
H_{\ell\ell'}(A):\mathfrak{su}(3)\to\mathfrak{su}(3).
\]
A standard block-matrix estimate gives
\[
\|\nabla^2 S_W(A)\|_{\mathrm{op}}
\le \sup_{\ell\in\mathcal{E}} \sum_{\ell'\in\mathcal{E}} \|H_{\ell\ell'}(A)\|_{\mathrm{op}}.
\]
Fix a link $\ell$. Only plaquettes $p\ni \ell$ contribute to the row $\ell$, and each such plaquette couples $\ell$ to at most three other links plus itself. By Lemma~\ref{lem:plaquette-blocks}, the row-sum contribution of one plaquette is at most
\[
\frac{\beta}{\sqrt{3}} + 3\cdot \frac{\beta}{3} = \beta\left(1+\frac{1}{\sqrt{3}}\right).
\]
Since at most $n_p=6$ plaquettes contain $\ell$, the row sum is bounded by $n_p(1+1/\sqrt{3})\beta$. Taking the supremum over $\ell$ proves the bound.
The looser estimate $8\sqrt{3}\beta$ follows from the inequality $6(1+1/\sqrt{3})\le 8\sqrt{3}$.
\end{proof}

\begin{remark}[Generalization and what this buys you]
The same argument works for $SU(N)$ with $\sqrt{3}$ replaced by $\sqrt{N}$, and $n_p$ replaced by the maximum plaquette incidence for the chosen lattice dimension.

This bound is \emph{uniform in the field amplitude} and \emph{uniform in the lattice volume}.
It is exactly the kind of analytic lemma that makes ``grid+Lipschitz+residual'' convexity certification realistic: one can combine it with a quadratic mass term (Haar or surrogate) to obtain global strong-coupling convexity, and locally one can bound third derivatives to extend verified eigenvalue checks from grid points to a whole ball.
\end{remark}

\end{document}
